\documentclass[11pt]{article}
\usepackage[top=2cm]{geometry}

\title{Homework 8: Keenan Crane's Biography}
\author{Alejandro Cornejo Pérez}
\date{March 2026}
\setlength{\parskip}{\baselineskip}
\begin{document}

\maketitle
Keenan Crane was born in 1982 in the United States and he is currently about 43 years old. He is a very well known researcher in the areas of computer graphics and geometric computing. his work has been focused on the mathematical and computational foundations of geometry processing. He has helped computers represent, analyze and manipulate complex shapes and surfaces. He has had a strong impact on areas like 3D modeling, animation, simulation, and scientific computing.

Crane studied computer science and mathematics during his undergraduate studies and later focused his career in computer science. He completed his PhD at the California Institute of Technology (Caltech). This is where he worked on developing mathematical methods for computing geometry on curved surfaces. His doctoral research helped the community by introducing new approaches for computing geometric quantities. Such as curvature, vector fields and differential operators.

One of the most important contributions made by Crane has been his work on discrete differential geometry. He adapted the tools of classical differential geometry so that they could be applied to digital surfaces used in computer graphics. These techniques allowed computers to perform accurate geometric computations on triangle meshes. Triangle meshes are the standard way of representing 3D objects in computer graphics.

He has also contributed by creating algorithms for geometry processing like mesh parameterization, surface flows, vector field design, and geometric optimization. These algorithms are very used in graphics research and they have influenced software tools for modeling, animation, and simulation. His work helped connect mathematics and computer graphics.

After he completed his PhD, Crane joined the Carnegie Mellon University, where he became a professor in the Computer Science Department. He currently leads the research in geometry processing and computer graphics in this college. He also teaches courses that have to do with geometry and graphics. His educational resources like lecture notes and open source software libraries, are used by students and researchers around the world. In conclusion, Keenan Crane has received several awards for his research contributions. He has played an important part in the advancing of mathematics in modern computer graphics. He created powerful tools for representing and manipulating complex 3D shapes.
\end{document}
